\section{Implementation Invariants}
\label{app:implementation_invariants}

The response is constructed first on the complete time-major, sensor-minor row
index.  The PM\(_{2.5}\) observation mask then selects the same ordered rows
from \(Y\), \(H_\Phi^{\mathrm{lag}}\), \(Q\), and metadata.  Source--basis
columns are always source-major and basis-minor.  Any row or column metadata
mismatch is a hard error.

The response construction uses integer source-cell centers, half-cell domain
extents, fractional final advection substeps, a positive minimum dispersion
age, direct sensor evaluation, finite lag, and complete removal after center
exit.  Kernel boundary/truncation loss and whole-release exit loss are recorded
separately.  Retained and dropped kernel mass sum to emitted kernel mass within
the recorded quadrature tolerance; sensor observations are never summed as a
mass proxy.  The default puff baseline has zero source and zero initial state.

Normal backgrounds have effective rank at most eight and do not use inventory
or response columns.  Thin-SVD projection stores the numerical tolerance,
singular values, retained \(U_r\), dependent input columns, removed components,
and orthogonality/idempotence residuals.  A source-like column is permitted only
when the run is labeled as a stress test.

The primary \(Q\), lag-candidate grid, and fixed-zero coefficient set are
declared before fitting.  A fixed-zero mask removes columns before diagnostics
and fitting and stores both reduced-to-original and original-to-reduced index
maps.  Fitted near-zero coefficients never modify the mask.  Adjacent lag
responses use identical rows and columns; the default convergence threshold is
\(\tau_L=10^{-3}\).  The reported runs use the declared \(L=12\); the lag sweep of
Experiment 7 records the largest candidate as its ``selected'' lag when no candidate
meets \(\tau_L\).

\section{Sanity Gates}
\label{app:sanity_gates}

Before paper experiments, the implementation must pass six small deterministic
gates built from public APIs:
\begin{enumerate}[leftmargin=*]
    \item \textbf{S1: response.}  On a \(16\times16\) grid with four compact
    sources, four sensors, impulse/constant bases, and seed 123, verify response
    shape, downwind arrival, open-boundary exit, no wraparound, dispersion
    anisotropy, mass accounting, metadata, and exact repeatability.
    \item \textbf{S2: projection.}  Use the S1 response with a rank-four normal
    background to verify exact ordering, background removal, orthogonality,
    idempotence, redundant-column invariance, and the labeled source-like stress
    behavior.
    \item \textbf{S3: diagnostics.}  Verify orthogonal, duplicate, and weak
    matrices, including padded zeros, \(\sigma_J\), infinite deficient
    condition status, weak-pair N/A values, and a matched multi-source eastward
    versus two-direction wind comparison.
    \item \textbf{S4: fit.}  Recover known nonnegative coefficients in
    noiseless/noisy well-conditioned cases, verify projected-FISTA KKT
    convergence, preserve exact nonnegativity, and recover the duplicate-pair
    sum when the individual split is nonunique.
    \item \textbf{S5: report groups.}  Verify the pairwise report-group rule used in
    the recorded controlled runs: the duplicate-source connected component, an
    \(A\)--\(B\)--\(C\) chain with both trigger edges retained, the separated-source
    non-edge, the independent weak flag, and the summed group contribution.
    \item \textbf{S6: end to end.}  Build sources, wind, response, background,
    masked observations, diagnostics, fit, uncertainty, and report groups with
    one command and write a compact provenance-complete summary.
\end{enumerate}
The experiments begin only after S1--S6 and the PyTorch CPU/CUDA parity
checks pass.  \textbf{S8 (grouping)} covers Theorem~\ref{thm:groupings}: the matroid components
equal the unique finest identifiable partition found by brute force on random
operators, a zero column is its own component, and the declared partition is
respected.  Two further gates support the results.  \textbf{S7 (adequacy
calibration)} runs a statistically calibrated study of the residual-adequacy test
on a controlled orthonormal design, checking that the null rejection rate matches
\(\alpha\), that the null \(p\)-values are uniform, and that power increases
monotonically with omission strength.  A final \textbf{reporting gate} exercises
the evaluation layer that emits the result tables in Appendix~\ref{sec:evaluation},
verifying our reporting rules: undefined weak-pair metrics stay null,
grouped metrics accompany any non-singleton report component, percentages are labeled fractions of fitted
sensor signal, and an uncalibrated run never reports an adequacy pass.

\section{PyTorch Solver and Uncertainty Protocol}
\label{app:pytorch_solver}

After CSV/NPZ ingestion, numerical computation uses PyTorch tensors on an
explicit device and dtype.  The nonnegative coefficient objective is
\begin{equation}
f(\mathbf c)=
\|\widetilde Y-\widetilde H_\Phi\mathbf c\|_2^2
+\lambda\|\mathbf c-\mathbf c_0\|_2^2,
\qquad \mathbf c\ge0.
\label{eq:appendix_fista_objective}
\end{equation}
Every reported run uses \(\lambda=0\), which is problem~\eqref{eq:nnls_estimator}.
Projected FISTA uses a Lipschitz step derived from \(\sigma_1\) or a recorded
power-iteration estimate, clamps every iterate to the nonnegative orthant, and
restarts acceleration whenever the objective increases.  Termination requires
the configured projected-gradient/KKT residual and relative objective-change
tolerances, or reports that the iteration cap was reached.  The fit artifact
stores convergence status, iteration count, final KKT residual, objective
summary, device, and dtype.

Active-set covariance uses PyTorch linear algebra on
\(\widetilde H_{\Phi,\mathcal A}^{\top}
\widetilde H_{\Phi,\mathcal A}+\lambda I\).  Wind-ensemble and bootstrap refits
are batched on the same device when shapes permit.  Undefined weak-pair metrics
are serialized as JSON \texttt{null}, never NaN or a false numerical distance.

Uncertainty artifacts distinguish \texttt{ensemble\_kind=transport} from
\texttt{ensemble\_kind=inventory}.  Only the former may be summarized as an
ensemble uncertainty interval; inventory members remain named robustness
scenarios, and an aggregation request that mixes the two kinds is an error.

\section{Adequacy and Ensemble Result Contracts}
\label{app:adequacy_contracts}

A residual-adequacy result stores \texttt{calibration\_status},
\(T_{\mathrm{res}}\), bootstrap quantile, Monte Carlo \(p\)-value,
\(\alpha\), replicate count, and noise-model provenance.  Paper runs use
\(\alpha=0.05\) and 1000 refitted replicates, except the structural-mismatch case of
Experiment 5, which uses 200.  If the externally calibrated
noise covariance is absent or invalid, numerical test fields are null and the
status is \texttt{uncalibrated}; no boolean pass is emitted.  Sensor, time, and
autocorrelation residual summaries are retained in either case.

A lag-selection result stores every candidate \(L\), adjacent convergence
ratio, selected \(L\), \(\tau_L\), physical grid rationale, row count, and
diagnostic/group summaries.  A wind-distribution result stores the sampled-window
provenance; 5th, 50th, and 95th percentiles of \(\sigma_J\); full numerical and
effective-rank probabilities; coefficient weakness probabilities; separation
distributions; and reported-partition frequencies.

\section{Reproducibility and Artifact Provenance}
\label{app:reproducibility}

Every paper run records the container image, repository revision, PyTorch and
CUDA versions, device model, dtype, deterministic-algorithm settings, random
seeds, source files and crop, per-source scales, wind-imputer architecture, held-out wind split, query grid, gridded wind product, and ensemble provenance,
inventory identifier/hash/version, observation mask, response and transport
configuration, wind realization, dispersion settings, background columns and
rank tolerance, lag protocol, fixed-zero mask and index mapping, diagnostic
thresholds, ensemble kind, solver settings, noise-model provenance, and parent
artifact hashes.
Generated wind products, response matrices, checkpoints, and experiment outputs
are written to generated-data or log paths and are not treated as immutable
source files.

Controlled-run artifacts contain true and estimated coefficients, reconstructed
activities, response/projection results, diagnostics, uncertainty, report
groups, and recovery metrics.  Observed New Delhi artifacts omit unavailable
ground truth and instead contain observed-row counts, fitted trajectories, raw
and projected residuals, normalized proxy activities, group signals,
uncertainty, weak/ambiguous flags, trigger pairs, and group summaries.  A
percentage is emitted only with an explicit denominator identifying it as a
fraction of fitted inventory-attributed sensor signal.

\noindent\textbf{Computing infrastructure.}
The paper-resolution controlled experiments (Experiments 1--10) and the observed fits
ran on one NVIDIA GPU with CUDA 11.8 and PyTorch 2.6; the response operator is stored in float32.  Experiments 1--10 were
re-run from the recorded configurations on one laptop CPU (Apple M5, PyTorch 2.14), and
every recorded scalar result agrees with the GPU run to within \(1\%\) relative
difference.
Experiment 11 and the signal-target re-analysis of the observed
weeks ran on that CPU; the re-analysis reads the committed projected operators and fits
and does not re-run the pipeline.

\noindent\textbf{Data and licences.}
The PM\(_{2.5}\) and meteorological records are public data of the Central Pollution
Control Board of India \citep{cpcb,cpcb_portal}.  The population grid is GPWv4,
Revision 11 \citep{ColumbiaUniversity2018}, distributed under the Creative Commons
Attribution 4.0 licence.  The brick-kiln and industry maps follow the regional
emissions inventory of \citet{GUTTIKUNDA2013101}, and the traffic map is derived from
road map data \citep{google_maps}; Appendix~\ref{app:platform_details} describes the
preparation.  The code is in the supplementary material.
